\scanpage{66}
\begin{center}6 - 34 -\hfill n\textsuperscript{o} 392\end{center}

6) Soient $k$ un corps, $B$ l'anneau de polynômes $k[X,Y]$, $A$ le sous-anneau $k[X^2,XY,Y^2]$ de $B$.

a) Montrer que $A$ est un anneau noethérien intégralement clos et que $B$ est un $A$-module de type fini.

b) Montrer que l'idéal $\underline p=AX^2+AXY$ de $A$ est un idéal premier de hauteur 1, (donc divisoriel), mais que le $B$-module $\underline p\otimes_A B$ a un module de torsion $\ne0$ et n'est donc pas réflexif ; l'idéal $\underline pB$ de $B$ n'est pas divisoriel, l'application canonique $\underline p\otimes_A B\longrightarrow\underline pB$ n'est pas injective et $\underline p$ n'est pas un $A$-module plat.

7) Soient $A$ un anneau noethérien intégralement clos, $E$ un $A$-module sans torsion de type fini, $E^*$ son dual.

a) Montrer que l'homomorphisme canonique $E^*\otimes_A E\longrightarrow\operatorname{End}_A(E)$ est un pseudo-isomorphisme.

b) Déduire de a) que pour que $E$ soit un $A$-module projectif, il faut et il suffit que $E^*\otimes_A E$ soit un $A$-module réflexif. (Remarquer que si $E^*\otimes_A E$ est réflexif, l'homomorphisme canonique $E^*\otimes_A E\longrightarrow\operatorname{End}_A(E)$ est bijectif).

*8) Soient $A$ un anneau noethérien intégralement clos, $M_1$, $M_2$ deux $A$-modules de type fini.

a) Montrer que les $A$-modules $\operatorname{Tor}_i^A(M_1,M_2)$, $\operatorname{Ext}_A^i(M_1,M_2)$ sont pseudo-nuls pour $i\geq2$, (se ramener au cas où $A$ est principal).

b) Si $M_1$ est sans torsion, montrer que $\operatorname{Tor}_1^A(M_1,M_2)$ et $\operatorname{Ext}_A^1(M_1,M_2)$ sont pseudo-nuls, (même méthode).

c) Si $M_1$ est un $A$-module de torsion, montrer que l'on a, (avec les notations du n\textsuperscript{o}\quad) :
\[
\begin{aligned}
\unsure{\chi|\gamma}(M_1\otimes_A M_2)-\unsure{\chi|\gamma}(\operatorname{Tor}_1^A(M_1,M_2))&=r(M_2)\unsure{\chi|\gamma}(M_1)\\
\unsure{\chi|\gamma}(\operatorname{Hom}_A(M_1,M_2))-\unsure{\chi|\gamma}(\operatorname{Ext}_A^1(M_1,M_2))&=-r(M_2)\unsure{\chi|\gamma}(M_1)\\
\unsure{\chi|\gamma}(\operatorname{Hom}_A(M_2,M_1))-\unsure{\chi|\gamma}(\operatorname{Ext}_A^1(M_2,M_1))&=r(M_2)\unsure{\chi|\gamma}(M_1)
\end{aligned}
\]

d) Montrer que pour tout couple de $A$-modules de type fini $M_1$, $M_2$, on a :
\[
\begin{aligned}
c(M_1\otimes_A M_2)-c(\operatorname{Tor}_1^A(M_1,M_2))&=r(M_1)c(M_2)+r(M_2)c(M_1)\\
c(\operatorname{Hom}_A(M_1,M_2))-c(\operatorname{Ext}_A^1(M_1,M_2))&=r(M_1)c(M_2)-r(M_2)c(M_1).\quad *
\end{aligned}
\]
\margin{65}
